% ----------------------------------------------------------------------
% Proposition: failure of the neural-network approximation on a general
% probability space. This file is an *artifact* cited in Section 4.2 of
% the thesis. It gives the counterexample showing that Buehler et al.
% (2019), Proposition 4.3, does NOT extend verbatim to general Omega in
% the standing Buehler setup (no (A1)--(A6) assumed): a one-period market
% with an integrable perfect hedge in which every network strategy has
% infinite entropic risk. The failure is located in (F3), not in (F1).
% When integrating into the final thesis, strip the standalone preamble
% and include this under the same numbering as the surrounding document
% (keep the "% draft: ..." comment to preserve the original label).
% Cross-references below are faked (hardcoded) so this artifact compiles
% without defining the targets; comments say what they should point to.
% ----------------------------------------------------------------------
\documentclass[11pt,a4paper]{report}
\usepackage[T1]{fontenc}
\usepackage[utf8]{inputenc}
\usepackage{amsmath,amssymb,amsthm}
\usepackage{enumitem}

\numberwithin{equation}{chapter}

\theoremstyle{plain}
\newtheorem{proposition}{Proposition}[chapter]
\newtheorem{lemma}[proposition]{Lemma}
\newtheorem{theorem}[proposition]{Theorem}
\newtheorem{corollary}[proposition]{Corollary}
\theoremstyle{definition}
\newtheorem{assumption}[proposition]{Assumption}
\newtheorem{definition}[proposition]{Definition}
\newtheorem{example}[proposition]{Example}
% Upright body, bold head (not the italic amsthm "remark" style).
\newtheorem{remark}[proposition]{Remark}

\newcommand{\E}{\mathbb{E}}
\newcommand{\PP}{\mathbb{P}}
\newcommand{\R}{\mathbb{R}}
\newcommand{\N}{\mathbb{N}}
\newcommand{\F}{\mathcal{F}}
\newcommand{\HH}{\mathcal{H}}
\newcommand{\Hu}{\mathcal{H}^{\mathrm{u}}}
\newcommand{\NNet}{\mathcal{N\!N}}
\newcommand{\as}{\ \PP\text{-a.s.}}

\begin{document}

\setcounter{chapter}{3}

\chapter{Artifact: failure on a general probability space}

\subsection*{The approximation property fails on general $\Omega$}

% FAKE REF: "Proposition~4.3" in the text below is the citation of the
%   result of Buehler et al.; in the thesis it refers to the reproduced
%   finite-space proposition (prop4_3_finite.tex, \ref{prop:finitespace}).
We now show that Proposition~4.3 of B\"uhler et al.\ does not extend
verbatim to a general probability space. The example is a one-period market with a perfect
(hence optimal) hedge whose hedge ratio is integrable but has no
exponential moments.

% draft: Proposition 4.7 (failure without extra assumptions);
% same construction as prop:counterexample in
% extension_of_approximation_on_general.tex.
\begin{proposition}\label{prop:counterexample}
  There exist a one-period market on a general probability space,
  without transaction costs and without trading constraints, an entropic
  risk measure $\rho_\lambda$ with arbitrary risk aversion $\lambda>0$
  and a position $X=-Z$ with $Z\in L^1$ such that
  \[
    \pi(X)\le0\qquad\text{and}\qquad\pi_M(X)=+\infty\quad\text{for every }M\in\N.
  \]
  Moreover, $\rho_\lambda(X+\delta_0(S_1-S_0))=+\infty$ for \emph{every}
  strategy of the form $\delta_0=F(I_0)$ with $F:\R\to\R$ bounded on
  $(0,1)$; this covers all networks with bounded activation and all
  networks with continuous activation.
\end{proposition}

\begin{proof}
  Let $\Omega:=(0,1)\times\{-1,+1\}$ with
  $\PP:=\mathrm{Leb}\otimes(\frac{1}{2}\delta_{-1}+\frac{1}{2}\delta_{+1})$,
  and write $U(\omega):=\omega_1$ and $\xi(\omega):=\omega_2$, so that
  $U$ is uniform on $(0,1)$, $\PP[\xi=\pm1]=\frac{1}{2}$ and $U,\xi$ are
  independent. Take one trading period ($n=1$), information $I_0:=U$,
  $I_1:=\xi$, hence $\F_0=\sigma(U)$ and $\F_1=\sigma(U,\xi)$, one asset
  ($d=1$) with
  \[
    S_0:=1,\qquad S_1:=1+\frac{\xi}{2},
  \]
  no costs ($c_k\equiv0$) and no constraints ($H_0(\delta_0):=\delta_0$).
  Note that $S$ is a $\PP$-martingale and
  $S_1>0$. Define
  \[
    h(u):=\frac{2}{\lambda}\log\frac{1}{u},\qquad u\in(0,1),\qquad
    Z:=h(U)\frac{\xi}{2},\qquad X:=-Z.
  \]
  Then $\E[|Z|]=\frac{1}{2}\E[h(U)]=\frac{1}{\lambda}<\infty$, so
  $Z\in L^1$.

  \textit{The perfect hedge.}
  The strategy $\delta_0:=h(U)$ is $\F_0$-measurable, and
  \[
    X+\delta_0(S_1-S_0)=-h(U)\frac{\xi}{2}+h(U)\frac{\xi}{2}=0,
  \]
  so $\pi(X)\le\rho_\lambda(0)=0$. Note that $\delta_0=h(U)\in L^1$; the
  optimal strategy is integrable, so Hornik's theorem applies to it and
  networks approximate it in $L^1(\PP)$ and, along a subsequence, almost
  surely.

  \textit{Every locally bounded strategy has infinite risk.}
  Let $\delta_0=F(U)$ with $m:=\sup_{u\in(0,1)}|F(u)|<\infty$. Then
  \[
    \rho_\lambda\bigl(X+F(U)(S_1-S_0)\bigr)
    =\frac{1}{\lambda}\log\E\Bigl[\exp\Bigl(\frac{\lambda}{2}\bigl(h(U)-F(U)\bigr)\xi\Bigr)\Bigr].
  \]
  The integrand is nonnegative, so the expectation is at least as large as
  its restriction to the event $\{\xi=+1\}\cap\{h(U)>m\}$. On that event one
  has $h(U)-F(U)\ge h(U)-m>0$, hence
  \[
    \exp\Bigl(\frac{\lambda}{2}\bigl(h(U)-F(U)\bigr)\xi\Bigr)
    =\exp\Bigl(\frac{\lambda}{2}\bigl(h(U)-F(U)\bigr)\Bigr)
    \ge\exp\Bigl(\frac{\lambda}{2}\bigl(h(U)-m\bigr)\Bigr)
    =\frac{e^{-\lambda m/2}}{U}.
  \]
  Independence of $U$ and $\xi$ and $\PP[\xi=+1]=\frac{1}{2}$ therefore yield,
  with $u_m:=e^{-\lambda m/2}$,
  \[
    \E\Bigl[\exp\Bigl(\frac{\lambda}{2}\bigl(h(U)-F(U)\bigr)\xi\Bigr)\Bigr]
    \ge\frac{e^{-\lambda m/2}}{2}\int_0^{u_m}\frac{1}{u}\,du=\infty,
  \]
  so the entropic risk is $+\infty$.

  \textit{Networks are locally bounded.}
  If the activation $\sigma$ is bounded, every network
  $F=W_L\circ F_{L-1}\circ\dots\circ F_1$ is bounded on all of $\R$,
  because the output of the last hidden layer lies in a bounded set and
  $W_L$ is affine. If $\sigma$ is continuousn, $F$ is continuous on $\R$ and therefore bounded on the
  compact set $[0,1]\supseteq(0,1)$. In either case, every network
  strategy $\delta_0=F(I_0)$ has $\sup_{(0,1)}|F|<\infty$, so
  $\rho_\lambda(X+W(\delta))=+\infty$ and $\pi_M(X)=+\infty$ for every
  $M$.
\end{proof}

\end{document}
